% ---------------------------------------------------------------------------
% Response to R1 Major 5 -- JEDC revision
% Body only: paste after the previous comment in the combined response file.
% Needs booktabs. Top heading is \section (numbered), subheadings \subsection*,
% matching the convention used for R2 Major 6.
% ---------------------------------------------------------------------------

\section{R1, Major 5}
{\itshape Lender-side risk pricing}

\begin{quote}\itshape
The precautionary-borrowing result assumes risk-neutral lenders. [\ldots] the
paper should discuss how robust it is to lender-side risk pricing, state that
its strength is conditional on risk-neutral pricing, and if feasible add a
robustness exercise with an exogenous SDF or a reduced-form global risk premium.
\end{quote}

We have run the exercise. The result does not depend on risk-neutral pricing:
at an empirically disciplined risk premium the state dependence is slightly
\emph{stronger} than in the baseline, and the debt ratio is essentially
unchanged.

\subsection*{The exercise}

The model carries an exogenous stochastic discount factor,
\[
  m' = \exp(-r^{*})\,\frac{\exp(\kappa\varepsilon')}{Z(p,\sigma')},
\]
where $\varepsilon'$ is the innovation to the log oil price and $Z$ normalises
the kernel on the discrete grid, so that $\mathrm{E}[m'] = \exp(-r^{*})$ holds
exactly at any $\kappa$. A positive $\kappa$ makes the discount factor high
when oil is dear, so lenders require more for default risk concentrated in those
states. Setting $\kappa = 0$ recovers the baseline exactly.

\subsection*{Disciplining $\kappa$}

$\kappa$ is not observable, so we discipline it by the maximum Sharpe ratio the
kernel implies, which for this kernel is $\sqrt{\exp(\kappa^{2}v)-1}$ in a state
whose conditional innovation variance is $v$. We average that across the
volatility states at their simulated frequencies and annualise. Equities run
about 0.4 a year, so a kernel implying much more than that is not a plausible
description of lenders.

\subsection*{The result}

\begin{table}[h]
\centering\footnotesize
\setlength{\tabcolsep}{5pt}
\begin{tabular}{ccccccc}
\toprule
$\kappa$ & mean spread & debt/GDP & default \% p.a. & \multicolumn{2}{c}{spread} & implied \\
\cmidrule(lr){5-6}
         &             &          &                 & low vol.\ & high vol.\ & Sharpe \\
\midrule
0    & 3.59\% & 26.1\% & 3.35\% & 3.27\% & 4.05\% & 0.00 \\
1.0  & 4.15\% & 26.1\% & 1.86\% & 3.79\% & 4.68\% & 0.28 \\
2.0  & 4.67\% & 25.1\% & 0.92\% & 4.25\% & 5.28\% & 0.56 \\
3.0  & 5.12\% & 23.3\% & 0.40\% & 4.67\% & 5.79\% & 0.87 \\
\bottomrule
\end{tabular}
\end{table}

\noindent
$\sigma_{\text{oil}} = 0.25$, every other parameter at its baseline value.
Spreads are whole-simulation means, debt is relative to annual output, and low
and high volatility are the extreme states of the simulated grid. The levels
are therefore not the Table 3 statistics, which average over the hundred
quarters preceding each default episode: on that basis the baseline
economy has a mean spread of 4.08\% and debt/GDP of 34.3\%, against targets of
4.1\% and 34\%. Every column is computed the same way at every $\kappa$, so the
comparison across rows is like for like.

\textbf{The state dependence strengthens monotonically in $\kappa$.} The gap
between high- and low-volatility spreads rises from $+0.77$ to $+1.12$. Risk
pricing loads on exactly the states where default risk is concentrated, so it
reinforces the mechanism rather than competing with it.

\textbf{At a plausible risk premium the calibration barely moves.} The equity
benchmark of 0.4 falls between the second and third rows of the table, at
$\kappa = 1.43$. There the gap is $+0.95$ against $+0.77$, about a fifth
larger, and debt/GDP is 25.6\% against 26.1\%. The precautionary-borrowing
result therefore does not rest on risk-neutral pricing, and the baseline is if
anything the conservative case.

\subsection*{What the exercise costs}

The default frequency falls sharply, from 3.35\% a year to 1.86\% at
$\kappa = 1$ and 0.92\% at $\kappa = 2$. That is the mechanical consequence of
pricing default risk more expensively: the country defaults less because
defaulting has become dearer to finance. The exercise holds every other
parameter at its baseline value, which isolates the risk-pricing channel but
means the $\kappa > 0$ economies no longer match the baseline's default
frequency. Re-matching it would require re-selecting $d_{1}$, which we have not
done, since the point is the comparative static rather than a second
calibration.

So the claim is about the sign and the ordering rather than the levels:
introducing a plausible risk premium leaves the state dependence intact and
slightly larger, at a debt ratio essentially unchanged, while lowering the
default rate below its target.

\subsection*{Suggested manuscript wording}

\begin{quote}
Spreads in the baseline are pure expected-default compensation, with no price of
risk. To check that the state dependence does not rest on this, we add an
exogenous stochastic discount factor in which the discount factor is high when
oil is dear, normalised on the solution grid so that the risk-free rate is
unchanged for any strength of risk pricing. The strength is disciplined by the
maximum Sharpe ratio the kernel implies: an annualised 0.4, the equity
benchmark, corresponds to $\kappa = 1.43$. At that value the gap between
high- and low-volatility spreads is 0.95 percentage points against 0.77 in the
baseline, and the debt-to-GDP ratio is 25.6 per cent against 26.1 per cent, so
risk pricing strengthens the state dependence slightly and leaves the debt
moment intact. The model's default frequency falls as risk pricing rises, from
3.35 to 1.86 per cent a year at $\kappa = 1$, so the risk-pricing economies no
longer match the default frequency of Table 3; we report the comparative static
rather than recalibrating the default cost to restore it.
\end{quote}
